\chapter{Results and Improvements}
\noindent\rule{\linewidth}{2pt}
\section{Introduction}

This chapter presents the results obtained from the implementation of the Airline Reservation System. It evaluates the performance of the system and discusses the improvements achieved over traditional manual methods. Additionally, it highlights possible enhancements that can be implemented in the future.

\section{Results}

The developed Airline Reservation System was successfully implemented using Java and tested with various user inputs. The system performs all core functionalities efficiently, including user registration, login, flight search, ticket booking, and cancellation.

\subsection{Functional Results}

The system was tested under different scenarios to verify its correctness and reliability. The following results were observed:

\begin{itemize}
\item The user registration module successfully stores user details without errors.
\item The login system correctly authenticates users based on valid credentials.
\item Flight details are displayed accurately with real-time seat availability.
\item Ticket booking is processed efficiently with proper seat allocation.
\item Cancellation of tickets updates the system records correctly.
\end{itemize}

\subsection{Performance Analysis}

The system demonstrates efficient performance in handling user operations. The response time for booking and cancellation operations is minimal, ensuring a smooth user experience. Since the system is console-based, it requires minimal system resources and performs efficiently even on low-end systems.

\subsection{Comparison with Existing System}

Compared to traditional manual reservation systems, the developed system offers several advantages:

\begin{itemize}
\item Reduces human errors in booking and record management.
\item Provides faster processing of reservations.
\item Ensures accurate and organized data storage.
\item Improves overall efficiency and user satisfaction.
\end{itemize}

\section{Improvements}

Although the current system meets basic requirements, several improvements can be made to enhance its functionality and usability.

\subsection{Graphical User Interface (GUI)}

The current system is console-based. A graphical user interface can be developed to make the system more interactive and user-friendly.

\subsection{Online Payment Integration}

The system can be enhanced by integrating secure online payment gateways to allow users to complete transactions digitally.

\subsection{Web-Based Application}

The system can be converted into a web-based application using technologies such as HTML, CSS, JavaScript, and frameworks like Spring Boot to allow access from anywhere.

\subsection{Real-Time Database Integration}

Currently, the system uses basic database operations. It can be improved by implementing real-time database updates and cloud-based storage for better scalability.

\subsection{Security Enhancements}

Additional security features such as encryption, authentication mechanisms, and role-based access control can be implemented to protect user data.

\section{Summary}

This chapter presented the results of the Airline Reservation System and demonstrated its effectiveness in handling booking operations. The system significantly improves efficiency compared to manual methods. Future enhancements can further expand the system’s capabilities and usability.